\documentclass[12pt]{article}
\usepackage[utf8]{inputenc}
\usepackage[T1]{fontenc}
\usepackage{hyperref}
\title{Pinigura Language Resources for K-12 Curriculum Development}
\author{Compiled by Claude Sonnet 5.5 (claude-sonnet-5-5) for Andr\'as Kornai}
\date{}
\begin{document}
\maketitle

\section{Introduction}
Pinigura (also Pinikura, Binigura; ISO 639-3 pnv) is a dormant Pama-Nyungan language of the Kanyara subgroup, formerly spoken along the Ashburton River in the Pilbara region of Western Australia. Ten speakers were recorded in 1975 and none since; a 2004 Wangka Maya survey found no full speaker, only a few people with partial or passive knowledge. Binigura people now speak Thalanyji. The earliest recordings are by O'Grady (1966), and Austin produced a wordlist.

\section{Primary Educational Resources}
No primers, textbooks or school programmes were found for Pinigura. Any heritage curriculum would likely build on the closely related Thalanyji.

\section{Linguistic and Reference Materials}
\begin{itemize}
\item Wikipedia, Pinikura language: \url{https://en.wikipedia.org/wiki/Pinikura_language}
\item Glottolog entry (pini1244): \url{https://glottolog.org/resource/languoid/id/pini1244.bigmap.html}
\item Wikipedia page reached through the ISO 639 code pnv, which covers the related Thalanyji language: \url{https://wikipedia.com/wiki/ISO_639:pnv}
\end{itemize}

\section{Religious and Community Resources}
\begin{itemize}
\item Wangka Maya Pilbara Aboriginal Language Centre, Binigura cultural introduction: \url{https://www.wangkamaya.org.au/pilbara-languages/binigura/01-cultural-introduction}
\end{itemize}

\section{Digital and Audio Resources}
\begin{itemize}
\item Endangered Languages Project entry: \url{https://endangeredlanguages.com/lang/4984?hl=en}
\end{itemize}
Archival recordings and the Austin wordlist are mentioned in the sources, but no links to them were found.

\section{Curriculum Development Context}
This is a revitalization situation with sparse documentation. The next steps would be to locate the O'Grady and Austin materials in archives (for example through the language centre or AIATSIS), consult Binigura and Thalanyji community members and Wangka Maya, and decide how much to borrow from Thalanyji. Nothing here should be presented as a teachable language course.
\end{document}